\honest{Scientific Evidence, Legal Evidence, Rational Evidence}{Eliezer Yudkowsky}{2007}

Your friend the police commissioner tells you in confidence that the city's crime boss is Wulky Wilkinsen. Should you believe it? Insulting Wulky would be foolhardy, I say, so the statement ``must have been strong Bayesian evidence.'' \nb{But insulting someone with even a small chance of being a crime boss is foolhardy. The same caution would follow from weak evidence.}

Courts will not jail Wulky on the commissioner's word. My explanation is that if we jailed everyone the commissioner named, corruption would grow, ``like honey attracts flies.'' \nb{American law gives its own reason. A witness must testify from personal knowledge, and testimony must be open to cross-examination, which tests the witness's perception and memory. These are rules about reliability. The post presents them as a social convention laid on top of rationality.}

So legal evidence is a subset of rational evidence. Then science. I am wearing white socks as I write this. You may believe me, and I could testify to it in court, but it is not science, because you cannot test it yourself. Science, I say, is made of generalizations that let you ``verify for yourself that the generalization is true, without having to trust anyone's authority.'' \nb{In practice almost no one, scientists included, checks more than a sliver of it.}

History, I say, is not science either, because our knowledge of Alexander the Great depends on authorities such as Plutarch. \nb{By this rule, what geology and cosmology say about events that happened only once, such as the extinction of the dinosaurs, is not scientific knowledge either.}

Is the prediction that the Sun will rise a month from now scientific? It may seem perverse to deny it, I admit, since it follows from generalizations you can test. Then I suggest that it is not: if you already knew the scientific belief about an experiment's result, ``why bother to run the experiment?'' \nb{By the same reasoning, no prediction would count as scientific before it had been tested.}

Is my definition of scientific knowledge true? ``That is not a well-formed question,'' I say; the standards are pragmatic choices. \nb{So the definition is stated as a fact until it is questioned, and then it becomes a convention.}

Perhaps, I suggest, future generations will count as scientific only papers published in open-access journals, and I say I think that would serve science better.

I end by saying that however we define science, a paper in an expensive closed journal and the commissioner's private word both still count as Bayesian evidence.
